\chapter{Glosario razonado}\label{ap:glosario}

Este glosario sirve para volver a un concepto después de haberlo estudiado. Los términos se
agrupan por las preguntas que ayudan a resolver. Las referencias llevan al capítulo donde
se desarrollan; una definición breve no sustituye sus supuestos ni sus ejemplos.

\section{Qué se observa}
\begin{description}[style=nextline,leftmargin=1.2em]
\item[Actividad de fuego detectada.] Propiedad del registro de un sistema de observación:
cuántas detecciones, presencias u otras unidades se obtuvieron bajo unos filtros. No equivale
automáticamente a actividad real total en el suelo. Capítulos~\ref{cap:problema} y~\ref{cap:calendario}.
\item[Anomalía térmica o foco de calor.] Píxel que satisface los criterios de un algoritmo
de detección térmica. La coordenada distribuida representa la geolocalización del píxel;
no identifica necesariamente el punto exacto de la llama. Capítulo~\ref{cap:problema}.
\item[Fuego activo.] Combustión que está ocurriendo durante la observación. Una pasada
puede perder una quema que empezó y terminó entre dos visitas. Capítulo~\ref{cap:orbitas}.
\item[Incendio.] Evento físico con continuidad espacial y temporal cuya delimitación
requiere reglas adicionales. No hay una correspondencia uno a uno entre incendios y filas
de FIRMS. Capítulo~\ref{cap:problema}.
\item[Área quemada.] Superficie afectada por combustión durante un intervalo. Se estima
con productos y métodos diferentes de los puntos de fuego activo; no se obtiene multiplicando
el número de focos por el área nominal del píxel. Capítulos~\ref{cap:problema} y~\ref{cap:frp}.
\item[Fuente térmica persistente.] Fuente de calor que aparece repetidamente, como una
instalación industrial o un volcán. Su recurrencia es una pista, no una clasificación
infalible. Capítulo~\ref{cap:errores}.
\item[Píxel y huella.] El píxel es la unidad de muestreo de la imagen; su huella es el
área del terreno a la que responde la medición. Su tamaño y forma varían con la geometría.
Capítulo~\ref{cap:escaneo}.
\item[Píxel mixto y fracción subpíxel.] Un píxel mezcla contribuciones de materiales o
temperaturas diferentes. La fracción subpíxel del fuego es la parte de su huella ocupada
por combustión, no la probabilidad de que haya un incendio. Capítulo~\ref{cap:deteccion}.
\end{description}

\section{Cómo llega la señal}
\begin{description}[style=nextline,leftmargin=1.2em]
\item[Longitud de onda.] Distancia entre dos puntos equivalentes de una onda; selecciona
una región del espectro. Un micrómetro es $10^{-6}$ metros. Capítulo~\ref{cap:radiacion}.
\item[Radiancia espectral.] Potencia por área proyectada, ángulo sólido e intervalo de
longitud de onda. Es una magnitud direccional; no debe confundirse con el flujo integrado
sobre un hemisferio. Capítulo~\ref{cap:radiacion}.
\item[Cuerpo negro.] Idealización que absorbe toda la radiación incidente y cuya emisión
térmica depende de la temperatura según Planck. Es una referencia física, no una
descripción literal de cada fuego. Capítulo~\ref{cap:radiacion}.
\item[Emisividad.] Cociente entre la emisión de un cuerpo y la de un cuerpo negro a la
misma temperatura, para una longitud de onda y dirección. Puede variar espectralmente.
Capítulo~\ref{cap:radiacion}.
\item[Temperatura de brillo.] Temperatura del cuerpo negro que produciría la radiancia
medida en un canal. Mezcla, atmósfera y emisividad impiden identificarla sin más con la
temperatura física de la llama. Capítulo~\ref{cap:radiacion}.
\item[Transmitancia.] Fracción de una señal que atraviesa un medio. Depende de la longitud
de onda y del camino recorrido, entre otros factores. Capítulo~\ref{cap:atmosfera}.
\item[Ventana atmosférica.] Región espectral de transmisión relativamente alta. No es
una garantía de visión a través de cualquier nube. Capítulo~\ref{cap:atmosfera}.
\item[SWIR, MIR y TIR.] Nombres de regiones del infrarrojo: onda corta, medio y térmico.
Sus límites dependen de la convención; al comparar sensores se citan las bandas y sus
longitudes concretas. Capítulos~\ref{cap:radiacion} y~\ref{cap:modisviirs}.
\item[Falso color.] Representación que asigna bandas, algunas invisibles al ojo, a los
canales rojo, verde y azul. Un color mostrado no es por sí mismo una temperatura ni una
clase de terreno. Capítulos~\ref{cap:problema} y~\ref{cap:atmosfera}.
\item[Saturación.] Estado en que la señal supera el rango utilizable del detector. Puede
permitir reconocer un píxel muy caliente pero impedir recuperar bien su radiancia.
Capítulo~\ref{cap:modisviirs}.
\item[FRP y FRE.] FRP es potencia radiativa del fuego, energía por unidad de tiempo,
normalmente expresada en MW. FRE es energía radiativa integrada en el tiempo. Sumar FRP
de pasadas separadas sin un intervalo temporal no produce energía. Capítulo~\ref{cap:frp}.
\end{description}

\section{Desde qué instrumento y en qué momento}
\begin{description}[style=nextline,leftmargin=1.2em]
\item[Sensor y plataforma.] El sensor es el instrumento; la plataforma es el satélite
que lo transporta. VIIRS en Suomi NPP y VIIRS en NOAA-20 comparten una familia instrumental,
pero representan oportunidades distintas de observación. Capítulo~\ref{cap:modisviirs}.
\item[MODIS y VIIRS.] Familias de instrumentos multiespectrales. Los productos de fuego
estudiados aquí usan resoluciones nominales diferentes. Es necesario especificar también
plataforma, producto y colección. Capítulo~\ref{cap:modisviirs}.
\item[Nadir.] Dirección hacia el suelo situada verticalmente bajo el satélite. La
resolución nominal suele referirse a esa geometría. Capítulo~\ref{cap:escaneo}.
\item[Franja o swath.] Banda de terreno observada durante una pasada. Una franja más
ancha puede reducir huecos espaciales sin eliminar nubes ni interrupciones de datos.
Capítulo~\ref{cap:orbitas}.
\item[Barrido y traza.] Direcciones transversal al movimiento y a lo largo de él.
\texttt{scan} y \texttt{track} describen dimensiones asociadas a esas direcciones.
Capítulo~\ref{cap:escaneo}.
\item[Efecto pajarita.] Solape entre líneas de barrido que crece hacia los bordes de
ciertos sensores. Puede multiplicar muestras de un mismo terreno. Capítulo~\ref{cap:escaneo}.
\item[Órbita heliosíncrona.] Órbita cuyo plano precesa de modo que se mantiene
aproximadamente la hora solar de cruce. No implica observación continua de un lugar.
Capítulo~\ref{cap:orbitas}.
\item[Deriva orbital.] Cambio progresivo de parámetros como altura u hora de paso cuando
no se mantienen sus valores nominales. Puede modificar la muestra del ciclo diario del
fuego. Capítulos~\ref{cap:orbitas} y~\ref{cap:armonizacion}.
\item[UTC, hora oficial y hora solar.] UTC es la referencia de tiempo usada en los
archivos; la hora oficial depende del huso legal; la hora solar depende de la posición
del Sol y de la longitud. La aproximación por longitud da hora solar media, no elimina
la ecuación del tiempo. Capítulo~\ref{cap:firms}.
\end{description}

\section{Qué producto se ha descargado}
\begin{description}[style=nextline,leftmargin=1.2em]
\item[FIRMS.] Sistema de NASA que distribuye información sobre fuego activo. La tabla
de puntos es un producto de acceso, no una máscara completa de observación del terreno.
Capítulo~\ref{cap:firms}.
\item[Colección, versión y reprocesamiento.] Una colección identifica una generación
del producto. Una etiqueta de versión puede codificar además revisiones de procesamiento.
Reprocesar aplica un procedimiento actualizado al archivo histórico; no añade una nueva
observación física. Capítulos~\ref{cap:firms} y~\ref{cap:calendario}.
\item[NRT y calidad estándar.] NRT prioriza la disponibilidad rápida; el procesamiento
estándar incorpora la cadena y los auxiliares correspondientes a su producto científico.
La sustitución de registros puede revisar valores ya descargados. Capítulo~\ref{cap:firms}.
\item[Confianza de detección.] Indicador del algoritmo sobre la calidad de una detección.
Las escalas MODIS y VIIRS no son probabilidades calibradas equivalentes.
Capítulo~\ref{cap:deteccion}.
\item[Omisión y comisión.] Omisión: fuego de referencia no detectado. Comisión: detección
que no corresponde al objetivo de referencia. Las tasas usan denominadores distintos.
Capítulo~\ref{cap:errores}.
\item[Cobertura y dato faltante.] Cobertura describe oportunidades de observación;
faltante describe información no disponible. Una lista vacía de detecciones no permite
reconstruir ninguna de las dos por sí sola. Capítulo~\ref{cap:calendario}.
\end{description}

\section{Cómo se comparan las series}
\begin{description}[style=nextline,leftmargin=1.2em]
\item[Celda-día.] Par formado por una celda fija y un día. Puede registrarse su presencia
de detección una vez por plataforma; esa definición reduce conteos múltiples, pero no
iguala sensibilidades. Capítulo~\ref{cap:armonizacion}.
\item[Armonización.] Conjunto de decisiones que busca hacer comparables indicadores
obtenidos con sistemas diferentes. No recupera la ubicación de fuegos no observados.
Capítulo~\ref{cap:armonizacion}.
\item[Solape.] Periodo en que existen observaciones de ambos sistemas. Permite estimar
relaciones, pero la coincidencia en años no garantiza coincidencia de pasadas y cobertura.
Capítulo~\ref{cap:armonizacion}.
\item[Estrato.] Grupo de datos definido por variables relevantes, como región y estación.
Puede tener su propia relación de calibración si dispone de información suficiente.
Capítulo~\ref{cap:armonizacion}.
\item[Sesgo.] Desviación sistemática respecto del objetivo o referencia. Aumentar la
muestra no lo elimina automáticamente. Capítulo~\ref{cap:errores}.
\item[Sobredispersión.] Varianza mayor que la media en un conteo. Advierte que un modelo
Poisson simple puede subestimar la variabilidad. Capítulo~\ref{cap:estadistica}.
\item[Bootstrap o remuestreo.] Construcción de réplicas de la muestra para estudiar la
variación de un estimador. Debe respetar las dependencias mediante unidades o bloques
adecuados. Capítulo~\ref{cap:estadistica}.
\item[Intervalo de confianza.] Intervalo producido por un procedimiento con una cobertura
de repetición definida bajo supuestos. No es el rango de variación entre años ni un
intervalo predictivo. Capítulos~\ref{cap:estadistica} y~\ref{cap:calendario}.
\item[Validación temporal y fuga de información.] La validación usa periodos reservados;
hay fuga si información de esos periodos interviene en ajustar o elegir el método que
se evalúa. Capítulo~\ref{cap:estadistica}.
\item[Anomalía.] Diferencia respecto de una referencia declarada. No expresa por sí sola
causalidad, significación estadística ni emergencia. Capítulo~\ref{cap:calendario}.
\item[Percentil, mediana y MAD.] El percentil sitúa un valor dentro de una distribución;
la mediana es su centro por orden; MAD resume distancias absolutas a la mediana. Es
necesario declarar las convenciones y tratar escalas nulas. Capítulo~\ref{cap:calendario}.
\item[MAE, MAPE y WAPE.] Medidas de desacuerdo absoluto o relativo. MAPE divide por cada
observación y falla con ceros; WAPE divide la suma de errores absolutos por el total
observado y pondera más la actividad grande. Capítulos~\ref{cap:estadistica} y~\ref{cap:aplicacion}.
\item[Procedencia y reproducibilidad.] Procedencia identifica el origen y las
transformaciones. Reproducibilidad permite recalcular el resultado con entradas y
reglas identificables; no demuestra por sí sola validez científica. Capítulo~\ref{cap:aplicacion}.
\end{description}
